\chapter{SYSTEM DESIGN AND RELATED TECHNOLOGIES}

\section{Mechanical Design}

This section covers the mechanical aspects of the robotic assistant's design, including chassis, mobility systems, and structural components.

\section{Electrical Design}

This section covers the electrical systems, power management, and hardware integration aspects of the robot.

\section{Software Design}

\subsection{Mobile Robot}

\subsubsection{ROS 2 (Robot Operating System 2)}

ROS 2 is an open-source robotic middleware framework that has become the de facto standard for robot software development. Designed from the ground up to overcome limitations of the original ROS, it provides comprehensive support for modern robotic systems, including multi-robot teams, resource-constrained embedded platforms, real-time control systems, and operation across unreliable networks. ROS 2 establishes well-defined design patterns and architectural conventions for structuring complex robotic software, supporting modularity, scalability, and maintainability in autonomous systems. It targets production-grade deployment environments where reliability and determinism are critical, combining distributed middleware capabilities with tools and libraries that streamline development of sophisticated control, perception, and communication subsystems. For this project, ROS 2 serves as the core integration middleware connecting all robot subsystems.

\subsubsection{ROS Navigation Stack}

The ROS Navigation Stack provides a comprehensive framework for autonomous mobile robot navigation, accepting odometry data and sensor stream inputs and outputting velocity commands to the robot's motor controller. This creates a clean abstraction that separates high-level path planning logic from low-level motor actuation. The Navigation Stack requires correct configuration of several components: ROS middleware access, a tf transform tree defining coordinate frame relationships between hardware components, standardised sensor data publication using ROS message types, and tuned costmap and planner parameters reflecting the robot's geometry and dynamics. This modular architecture enables robust navigation in diverse environments while supporting independent testing and development of individual components.

\subsubsection{AMCL (Adaptive Monte Carlo Localization)}

Adaptive Monte Carlo Localization (AMCL) is the primary localisation package used within the ROS ecosystem for autonomous navigation tasks. The algorithm employs a particle filter approach, maintaining a probabilistic distribution of candidate robot poses and iteratively refining this distribution as new sensor observations and odometry measurements become available. AMCL is specifically optimised for use with 2D laser rangefinders and pre-built occupancy grid maps, making it well-suited to structured indoor environments.

It is important to note that AMCL carries specific operational constraints. The algorithm assumes a fixed geometric relationship between the laser sensor and the robot base; any movement of the laser relative to the base frame violates this assumption and can degrade localisation accuracy significantly. Furthermore, AMCL requires continuous odometry input to propagate the particle filter over time. Odometry estimates may be derived from multiple complementary sources, including wheel encoders and inertial measurement units (IMUs), which together provide more robust position and orientation estimates than either source alone.

\subsection{Database Management}

\subsubsection{Retrieval-Augmented Generation (RAG)}

Retrieval-Augmented Generation (RAG) is a technique that enhances large language model output quality by incorporating an external knowledge retrieval step before response generation. Rather than relying solely on patterns in the model's training parameters, RAG queries a domain-specific knowledge base at inference time, retrieving contextually relevant information to augment the model's context.

RAG offers cost-effective extension of LLM capabilities to specialised domains without expensive model retraining. By grounding responses in retrieved, current information, RAG substantially reduces factual hallucinations and improves relevance and accuracy. For this project, RAG forms the backbone of the knowledge retrieval subsystem, enabling contextually appropriate responses to user queries about the operating environment and domain-specific topics.

\subsubsection{Vector Database}

A vector database stores data as dense numerical vectors (embeddings) rather than conventional row-column structures, enabling similarity-based retrieval over large unstructured data collections including text, images, and audio. Vector databases support approximate nearest-neighbour search rather than exact-match queries, enabling retrieval of semantically related records even with differing query wording. This functionality is essential for RAG systems where semantic relevance rather than lexical overlap determines retrieval quality.

\subsubsection{Embeddings}

Embeddings are dense, fixed-dimensional numerical representations of objects — including text, images, and audio — produced by machine learning models trained to capture semantic and contextual relationships. Objects that are semantically similar are mapped to nearby points in the embedding space, enabling distance-based similarity measures to serve as a proxy for meaning. In the context of this system, text embeddings are used to represent both the stored knowledge base documents and incoming user queries, enabling the vector database to identify and retrieve the most relevant passages for augmenting the LLM's response.

\subsubsection{Milvus}

Milvus is an open-source vector database platform specifically engineered for large-scale similarity search on massive embedding collections. Designed with production deployment in mind, Milvus provides highly optimised indexing structures and query algorithms that enable fast, accurate retrieval of semantically similar vectors from datasets containing millions or billions of embeddings. Unlike traditional relational databases, Milvus is purpose-built for vector search operations, supporting multiple similarity metrics including Euclidean distance, cosine similarity, and inner product.

Key features of Milvus include exceptional scalability through distributed architectures that allow horizontal expansion, high-performance approximate nearest-neighbour search capabilities that retrieve relevant results with minimal latency even on very large datasets, and flexible indexing strategies optimised for different trade-offs between search accuracy and computational efficiency. In this project, Milvus serves as the underlying vector database infrastructure for the RAG pipeline, enabling rapid semantic retrieval of domain-relevant information from the knowledge base to augment LLM responses.

\subsubsection{SQL (Structured Query Language)}

SQL is a standardised declarative programming language for defining, manipulating, and retrieving data in relational database management systems (RDBMS). Data is organised in tables with rows representing records and columns representing typed attributes. SQL provides comprehensive operations for schema creation, record insertion and updates, referential integrity enforcement, and complex multi-table queries through joins and aggregations. Query optimisation and index management facilities ensure retrieval performance at scale. In this project, SQL maintains structured records of user interactions, profiles, and system logs, complementing the vector database layer for semantic retrieval.

\subsection{Artificial Intelligence}

\subsubsection{Face Recognition}

Facial recognition is a biometric identification technology that authenticates or identifies individuals by analysing distinctive geometric and photometric features of the human face. It detects facial regions, generates compact feature representations, and compares them against stored facial templates. The system automatically identifies returning users and personalises interactions based on stored preferences and interaction history. This capability enhances human-robot interaction naturalness and continuity by eliminating explicit identification steps. For the robotic assistant, facial recognition enables context-aware engagement with individual users.

\subsubsection{Natural Language Processing (NLP)}

Natural Language Processing (NLP) is a multidisciplinary field enabling computational systems to understand, interpret, and generate human language in written and spoken forms. It combines rule-based linguistics, statistical models, and deep learning architectures for tasks including parsing, sentiment analysis, intent detection, and machine translation. NLP systems capture not only surface language form but also deeper semantic content, including speaker intent, implied meaning, and emotional tone. This nuanced understanding enables the robotic assistant to engage in natural, contextually appropriate dialogue, especially in open-domain settings with ambiguous or colloquial queries.

\subsubsection{Large Language Model (LLM)}

A Large Language Model (LLM) is a class of artificial intelligence system trained on massive corpora of text data — often comprising hundreds of billions of tokens — to acquire broad and deep competencies in understanding and generating human language. LLMs are built on the transformer neural network architecture, which leverages self-attention mechanisms to capture long-range dependencies between tokens and model the statistical structure of language at multiple levels of abstraction.

Through large-scale pre-training followed by task-specific fine-tuning or instruction alignment, LLMs develop the capacity to perform a diverse array of language tasks — including question answering, summarisation, translation, and code generation — without requiring explicit task-specific training data for each application. In this project, the LLM serves as the central reasoning and response generation component of the conversational AI subsystem, processing retrieved contextual information from the RAG pipeline to produce coherent, factually grounded responses to user queries.

\subsubsection{Text-to-Speech (TTS)}

Text-to-Speech (TTS) is a speech synthesis technology that converts written text into natural-sounding spoken audio, enabling AI systems to communicate verbally with users without requiring pre-recorded audio assets. Modern TTS systems are built on deep learning architectures — including neural vocoders and sequence-to-sequence models — that generate speech with dynamic, context-sensitive variations in prosody, pitch, speaking rate, and pronunciation. These characteristics produce output that closely approximates natural human speech patterns, significantly improving comprehensibility and user engagement compared with earlier rule-based synthesis methods.

TTS technology has found widespread application in accessibility tools for visually impaired users, virtual assistants, and interactive voice response systems. In the context of this robotic assistant, TTS serves as the primary output modality for verbal communication, enabling the robot to deliver responses audibly in environments where visual screen interaction is impractical or undesirable.

\subsubsection{Automatic Speech Recognition (ASR)}

Automatic Speech Recognition (ASR) is a technology that processes raw acoustic speech signals and converts them into machine-readable text, forming the primary input pathway for voice-driven human-machine interaction. Contemporary ASR systems employ end-to-end deep learning architectures — including recurrent neural networks and transformer-based models — trained on large-scale speech corpora to achieve high transcription accuracy across a range of speakers, accents, and acoustic conditions.

Advanced ASR implementations integrate Natural Language Processing modules downstream of the transcription stage, enabling the system to infer the user's intent and semantic meaning from the recognised text in addition to producing a verbatim transcript. ASR accuracy is sensitive to several environmental factors, including background noise levels, microphone quality, and speaker proximity; consequently, robust noise handling and signal preprocessing are important design considerations for deployment in real-world robotic environments. Despite these challenges, modern ASR systems are capable of achieving near real-time transcription latency, enabling fluid conversational interactions between the user and the robotic assistant.